\section{Einführung}
\slide{Kognition}{
    \begin{itemize}
        \item \b{Enge Definition:} Alle bewussten Denkprozesse wie Entscheidungen, Abwägungen und Problemlösung.
        \item \b{Weite Definition:} Alle bewussten und unbewussten Denkprozesse.
        \item Kognition ist die Umgestaltung von Informationen durch ein verhaltenssteuerndes System.
    \end{itemize}
}
\slide{Kognitive vs. Nicht-Kognitive Systeme}{
    \begin{itemize}
        \item \b{Nicht-kognitive Prozesse:} Starre Kopplung von Reizen und Reaktionen, die nicht durch Lernen veränderbar ist.
        \item \b{Kognitive Prozesse:} Flexible Kopplung von Reizen und Reaktionen; Reaktionen sind durch Lernfähigkeit veränderbar und auswählbar.
    \end{itemize}
}
\slide{Grundhypothese der Kognitionswissenschaft}{
    \begin{itemize}
        \item Kognition wird als Informationsverarbeitung verstanden, die berechnet werden kann.
        \item \b{Zentrale Gleichung:} Kognition = Repräsentation + Prozess.
    \end{itemize}
}
\slide{Repräsentation und Prozess}{
    \begin{itemize}
        \item \b{Prozess:} Ein algorithmisch beschreibbarer Vorgang, der auf eine Repräsentation angewandt wird.
        \item \b{Verhältnis:} Repräsentation und Prozess sind voneinander abhängig und treten nicht getrennt auf.
        \item \b{Wissensrepräsentation:} Format von Informationen zur Enkodierung (z. B. im Computer).
        \item \b{Mentales Modell:} Ein unvollständiges Modell der Welt, das subjektiv, erfahrungsbasiert und kontextabhängig ist.
        \item \b{Abweichung:} Mentale Repräsentationen weichen oft von der objektiven Wirklichkeit ab (Beispiel: Fisch stellt sich Kuh vor).
    \end{itemize}
}
\slide{Behaviorismus vs. Kognitionswissenschaft}{
    \begin{itemize}
        \item \b{Behaviorismus:} Erklärt Verhalten rein durch Reiz-Reaktions-Abfolgen (Blackbox bleibt geschlossen).
        \item \b{Kognitionswissenschaft:} Untersucht die kognitiven Prozesse (die "Blackbox"), welche die Reaktionen zwischen Reiz und Verhalten beeinflussen.
    \end{itemize}
}
\slide{Methodik}{
    \begin{itemize}
        \item \b{Methoden:} Formale Analyse, empirische Forschung und Computer-Modellierung.
        \item \b{Methodische Iteration:} Ein zirkulärer Forschungsprozess, bei dem sich Theorie, Empirie und Modellierung gegenseitig beeinflussen und weiterentwickeln.
        \item \b{Ziel der Modellierung:} Menschliche Denkprozesse in Computern darstellen, um Rückschlüsse auf die menschliche Kognition zu ziehen.
    \end{itemize}
}





